\subsection{The marked-cusp inertia locus and its surviving central Frobenius classes}

\begin{statement}{E.67}{Proved}{Complete marked-cusp inertia classes and an explicit ramified class without a full lift}
Retain the actual cubic curve, the good place $v\nmid3$ of
$K=\mathbb Q(\zeta_3)$, and the local arithmetic group of E.61.
Write $q=|k_v|$ and let $Q=B_0$ be the single rational cusp
marked in E.41. Its actual positive peripheral generator is
$d=d_Q$. Set
\[
 I=\overline{\langle\!\langle d\rangle\!\rangle}_{\Pi},\qquad
 W=\Pi/\gamma_3I,\qquad B=\Pi/I,\qquad
 K_I=I/\gamma_3I.
 \tag{E67.1}
\]
Here $\Pi$ is the free pro-three group of rank ten and every
subgroup and commutator image is closed. Thus $W$ is exactly
the single-cusp marked-depth quotient in E41.1; the kernels
of $\Pi\to W\to B$ are respectively $\gamma_3I$ and $K_I$.
The quotient $B$ is the actual pro-three group of the cubic
with this one cusp filled, and is free of rank nine.

Use the actual cusp Frobenius $F_Q=z_QF$ of E.61, so that
$F_Q(d)=d^q$. Replacing the generator $F$ by $F_Q$ is the
literal change of arithmetic coordinates
$(w,F_Q)\mapsto(wz_Q,F)$; it retains the given arithmetic
extension and the same inertia generator. Denote the induced
automorphism of $B$ by $\theta$. The two actual normal-kernel
modules and their actions are
\[
 \begin{split}
 D&=I/I'=\mathbb Z_3[[B]],\\
 C&=\gamma_2I/\gamma_3I
       =\widehat{\bigwedge}^{2}_{\mathbb Z_3}D,\\
 F_Q|_D&=q\theta_*,\qquad
 F_Q|_C=q^2\theta_2,\qquad
 \theta_2=\widehat{\bigwedge}^{2}\theta_* .
 \end{split}
 \tag{E67.2}
\]
The element $bdb^{-1}$ maps to $\delta_b\in D$.
The completed exterior module in this formula means the
compact inverse limit
\[
 C=\varprojlim_U
       \bigwedge^2_{\mathbb Z_3}\mathbb Z_3[B/U],
 \qquad
 \mathcal Z_Q=C/(q^2\theta_2-1)C,
 \tag{E67.3}
\]
where $U$ runs through characteristic open subgroups of $B$.
All coefficient maps are the actual trace maps on the
displayed finite coset sets; the denominator is its compact
image and hence is already closed.

Consider the precise locus of local $W$-sections whose
geometric inertia becomes the identity in $B$. Its ramified
part, up to full geometric conjugacy in $W$, is
\[
 (\mathbb Z_3\setminus\{0\})\times
       \mathbb Z_3/(q-1)\mathbb Z_3\times\mathcal Z_Q.
 \tag{E67.4}
\]
For $(a,\bar b,\bar w)$ choose any $w\in C$ over $\bar w$.
The representative is the actual section
\[
       t\longmapsto d^at,\qquad
       F\longmapsto d^b c_wF_Q,
 \tag{E67.5}
\]
where $c_w$ denotes the element of the central subgroup
$C\subset K_I$. The parametrization is a homeomorphism.
Its complete equivalence relation is
\[
 b\longmapsto b+(1-q)k,\qquad
 w\longmapsto w+(1-q^2\theta_2)v
       \quad(k\in\mathbb Z_3,\ v\in C).
 \tag{E67.6}
\]
The remaining part of the specified locus has trivial
geometric inertia. It consists exactly of the pairs
$(1,z)$ for arbitrary $z\in W$, modulo
$z\mapsto g zF_Q(g)^{-1}$. Every such unramified class has
a full $\Pi$-lift obtained by lifting $z$.
There is no nontrivial inertia with zero image in $D$ in
this locus.

The actual full-to-marked map on the ramified locus
whose inertia dies in $B$ has image precisely
\[
 (\mathbb Z_3\setminus\{0\})\times
       \mathbb Z_3/(q-1)\mathbb Z_3\times\{0\}.
 \tag{E67.7}
\]
It is a homeomorphism onto this image, and each of these
classes has exactly one full section class above it.
This assertion specifies the selected locus; it does not
identify the inertia classes whose image in $B$ is nontrivial.

At $K_v=\mathbb Q_2(\zeta_3)$, so $q=4$, the finite
arithmetic quotient of $\mathcal Z_Q$ supplied by the
first moment map is
\[
 \mathcal Z_Q\twoheadrightarrow
   B^{\rm ab}/(16F_B-1)B^{\rm ab}
       \simeq(\mathbb Z/9)^5\oplus(\mathbb Z/27)^2.
 \tag{E67.8}
\]
Here $F_B$ is the actual quotient of the E.37 Frobenius
matrix by the cusp vector $v_Q$. In particular set
$b_0=d_{A_0}$ and use the literal word
\[
          c_*=[d,b_0db_0^{-1}]\in\gamma_2I/\gamma_3I.
 \tag{E67.9}
\]
Then
\[
       t\longmapsto dt,\qquad F\longmapsto c_*F_Q
 \tag{E67.10}
\]
is a ramified local $W$-section with no full $\Pi$-lift.
Its abelian inertia is the single primitive vector $v_Q$.
The finite moment character and the integral matrix
certificate below detect its nonzero third coordinate
in (E67.4).

\emph{Manuscript: \S51.16, folios733v--736v, scans172--178,
formulas(371)--(379), Part E(c),(d); the intrinsic single
marked cusp and its exact kernel are those of E.25 and
E.41.}
\end{statement}

\begin{proof}
\emph{The actual normal-kernel model.}
Complete $d$ to a free pro-three basis of $\Pi$. Its other
nine generators give a geometric splitting of
$\Pi\twoheadrightarrow B$. This splitting is only used
to label conjugates and is not asserted to be arithmetic.
The normal closure $I$ is the free pro-three group on the
compact generator space
\[
                 B\longrightarrow I,
                    b\longmapsto bdb^{-1}.
 \tag{E67.11}
\]
One can prove this universal description directly. Form the
free pro-three group on the profinite space $B$, with its
continuous $B$-action by left translation. Its semidirect
product with $B$ is generated by the generator at $1$ and
the chosen nine generators of $B$. The universal property
of $\Pi$ gives a map to this semidirect product. Conversely
(E67.11) gives a $B$-equivariant map back to $I$, hence to
$\Pi$. The two composites fix every indicated generator,
so are inverse. The free group on the profinite space is
constructed as the inverse limit of the free pro-three
groups on its finite clopen partitions; this also verifies
continuity of both maps.

At nilpotence class two, this gives the explicit group
\[
 K_I=D\times C,\qquad
 (h,w)(h',w')=
     (h+h',w+w'+\tfrac12h\wedge h').
 \tag{E67.12}
\]
The generator at $b$ is $(\delta_b,0)$, the commutator
of $(h,0)$ and $(h',0)$ is $(0,h\wedge h')$, and the
central group is $C$. Division by two is integral over
$\mathbb Z_3$. For each finite coset set $B/U$, this
formula is the free class-two pro-three group on that
set: collection gives its unique coordinates and proves
the universal property. Taking the compact inverse limit
proves (E67.12) with kernel exactly $\gamma_3I$.
Equivalently $C$ is the alternating part of
$\mathbb Z_3[[B\times B]]$: the map sends
$\delta_g\wedge\delta_h$ to
$\delta_{(g,h)}-\delta_{(h,g)}$, and the inverse divides
by two after the natural exterior projection. This
describes its topology, including the diagonal, exactly.

Since $F_Q(d)=d^q$, applying $F_Q$ to $bdb^{-1}$
gives $q\delta_{\theta(b)}$ in $D$. Its unknown inner
part in $I$ has no effect on that abelianization. Applying
commutators then gives $q^2\theta_2$ on $C$. Thus
(E67.2) follows from the actual action, with no choice of
an arithmetic splitting of $W\to B$.

\emph{Two integral measure facts for this Frobenius.}
The group $B$ is the good-reduction group of the curve
with $Q$ filled. Its abelianization has seven cusp Tate
directions and the rank-two compact elliptic quotient.
Frobenius eigenvalues have complex absolute values $q$
and $\sqrt q$, respectively, by the same actual filling
and specialization maps as in E.61. On every lower-central
graded piece of the free pro-three group $B$, all
eigenvalues are products of one or more of these
eigenvalues and have absolute value greater than one.
The integral free Lie comparison and separation of the
lower-central filtration follow by the Magnus argument
used in E.41, applied to the nine free generators.
Consequently
\[
                 B^{\theta^m}=1\qquad(m\ge1).
 \tag{E67.13}
\]
An element fixed by $\theta^m$ would have a lowest
nonzero graded coordinate fixed by that matrix, which
the eigenvalue calculation excludes.

For $z_B\in B$, consider the affine permutation
$T(b)=z_B\theta(b)$. Each power of $T$ has at most
one fixed point: the quotient of two such points would
be fixed by a power of $\theta$. A periodic point of
$T$ is therefore a fixed point. The closure of $T$ is
pro-three, because both the actual Frobenius action and
the group of translations have this property by E.61.
On characteristic finite quotients of $B$, orbit lengths
are consequently three-powers. Compactness, applied to
the fixed-point equations for each bounded power of $T$,
forces these lengths to grow uniformly away from its
possible fixed point. The trace argument of E66.16 gives
\[
 D^T=\begin{cases}
       \mathbb Z_3\delta_{b_*},&T(b_*)=b_*,\\
       0,&\text{there is no such point}.
      \end{cases}
 \tag{E67.14}
\]

We also need the finite-coefficient version with a scalar
twist. Set $D_0=\ker(D\xrightarrow{\rm aug}\mathbb Z_3)$;
then $D=\mathbb Z_3\delta_1\oplus D_0$ is
$\theta_*$-stable. For every nonzero $a\in\mathbb Z_3$,
\[
              (D_0/aD_0)^{\,q\theta_*=1}=0.
 \tag{E67.15}
\]
For a unit $a$ there is nothing to prove. Otherwise work
over $\mathbb Z/3^r$, $r=v_3(a)$, and let a measure
satisfy $q\theta_*\mu=\mu$. All nonidentity
$\theta$-orbits are infinite by (E67.13). On a refined
cycle its coefficients are related by powers of $q$.
The contribution to a fixed coarser cycle is a unit
multiple of a coefficient times the geometric sum
$\sum_{j=0}^{L-1}q^{j\ell}$, where $L$ is the ratio
of the cycle lengths. Since $q\equiv1\pmod3$, this
sum has three-adic valuation $v_3(L)$; this follows
by induction from
$1+x+x^2=3+3(x-1)+(x-1)^2$ for $x\equiv1\pmod3$.
Orbit-length growth therefore kills every coefficient
away from $1$ modulo $3^r$. The measure is supported
at $1$, and its zero augmentation makes it zero.
Finite-coefficient trace systems have surjective maps,
so they identify exactly with $D/aD$; the centered
splitting identifies the stated submodule. This proves
(E67.15) with the actual compact topology.

\emph{The inertia element, including its zero case.}
Use the arithmetic coordinates based on $F_Q$. A section
in the specified locus has $\nu\in K_I$ and
$z\in W$, with
\[
                     zF_Q(\nu)z^{-1}=\nu^q.
 \tag{E67.16}
\]
Its image $h\in D$ is fixed by the affine permutation
$T(b)=z_B\theta(b)$, since the action on $D$ is $qT$.
If $h=0$, then $\nu\in C$. The action on $C$ is
$q^2\bigwedge^2T$, so (E67.16) would make
$(q\bigwedge^2T-1)\nu=0$. On each finite coset
quotient the exterior permutation has finite order;
$q^n-1\ne0$ in $\mathbb Z_3$ for every $n>0$.
Thus this operator is injective at each finite integral
level, and hence on $C$. It follows that $\nu=1$.
This proves the full zero-inertia assertion, rather than
only a statement about its abelian image.

For $h\ne0$, (E67.14) gives $h=a\delta_{b_*}$ for
a unique $a\ne0$ and a fixed point $b_*$. Conjugating
by a lift of $b_*^{-1}$ sends $z_B$ to $1$ and $h$
to $a\delta_1$. We now have
\[
                \nu=d^a c_w,\qquad z=(h_z,w_z)\in K_I.
 \tag{E67.17}
\]
The genuine action fixes $d$ by its $q$-power and acts
on $C$ by $q^2\theta_2$. Formula (E67.12) turns the
relation into the exact central equation
\[
          (q\theta_2-1)w=a\,\delta_1\wedge h_z.
 \tag{E67.18}
\]
Decompose $C=(\delta_1\wedge D_0)\oplus
\widehat{\bigwedge}^2D_0$. On the second summand,
$q\theta_2-1$ is injective by the same finite-permutation
argument. Hence $w=\delta_1\wedge r$ for a unique
$r\in D_0$. Equation (E67.18) becomes
\[
       (q\theta_*-1)r
          =a\bigl(h_z-\operatorname{aug}(h_z)\delta_1\bigr).
 \tag{E67.19}
\]
By (E67.15), $r=av$ with $v\in D_0$.
Conjugating by the actual element $(v,0)\in K_I$
subtracts $a\delta_1\wedge v$ from the central
coordinate of $\nu$, so sends it to $d^a$.
Its effect on the first coordinate of $z$ is
$h_z+(1-q\theta_*)v
=\operatorname{aug}(h_z)\delta_1$.
The remaining central coordinate is retained as an
arbitrary element of $C$; no unknown higher action has
been dropped from that conjugation.

\emph{Centralizers, representatives and topology.}
In $W$, conjugation of the first coordinate $a\delta_1$
by a nonidentity element of $B$ moves its support.
Thus an element centralizing $d^a$ lies in $K_I$.
There (E67.12) says that its first coordinate $h$ must
satisfy $a\delta_1\wedge h=0$.
The integral exterior module is torsion-free and its
centered decomposition makes this equivalent to
$h\in\mathbb Z_3\delta_1$. Consequently
\[
                   C_W(d^a)=\langle d\rangle\times C
                     \qquad(a\ne0).
 \tag{E67.20}
\]
After inertia has been made $d^a$, the ratio of its
Frobenius lift and $F_Q$ centralizes that inertia, so
is exactly $d^bc_w$. Conversely every such element
gives (E67.5), because $C$ commutes with $d$.
All geometric equivalences are in the same centralizer
and give precisely (E67.6). Different nonzero $a$ have
different abelian inertia. The compact annuli
$v_3(a)=r$, together with the compact groups in the
other two coordinates, map continuously and bijectively
to the corresponding Hausdorff orbit spaces. They
therefore give homeomorphisms, and these clopen strata
exhaust the ramified locus.

By E.61, a full ramified section is at one of the nine
cusps. The other eight cusp vectors remain nonzero in
$B^{\rm ab}$, so only $Q$ occurs in the selected locus.
Its full representatives are $d^bF_Q$ with no central
factor. This proves (E67.7), including uniqueness of
the full class. For $\nu=1$, every lift of $z$ satisfies
the full tame relation; this proves the stated
unramified lifting map.

\emph{The finite moment quotient and the explicit failed lift.}
There is a continuous, surjective homomorphism
\[
 \begin{split}
 \mathfrak m_1:C&\longrightarrow B^{\rm ab},\\
 \delta_g\wedge\delta_h&\longmapsto[h]-[g].
 \end{split}
 \tag{E67.21}
\]
Under the alternating-measure identification following
(E67.12), it is one half of the integral of the
continuous function $(g,h)\mapsto[h]-[g]$.
Equivalently apply the first moment to
$\operatorname{aug}(x)y-\operatorname{aug}(y)x$.
These formulas agree on the dense finite sums and show
continuity at every finite coefficient level.
It is equivariant for simultaneous left translation,
which acts trivially on its target, and for $\theta$.
Thus its exact Frobenius action is $q^2F_B$ and it
induces the surjection in (E67.8).
For any chosen free basis $b_1,\ldots,b_9$ of $B$,
the vectors $\delta_1\wedge\delta_{b_j}$ give an
integral right inverse on its abelian basis.

At the place above two, the E.37 basis is ordered
$f_1,f_2,e_{A_0},e_{A_1},e_{A_2},e_{B_0},e_{B_1},
e_{B_2},e_{C_1},e_{C_2}$. The marked quotient removes
the sixth vector. In the remaining order its actual
matrix is
\[
 F_B=\begin{pmatrix}
 -2&0&0&0&0&0&0&0&0\\
 0&-2&0&0&0&0&0&0&0\\
 0&2&4&0&0&0&0&0&0\\
 -4&4&0&4&0&0&0&0&0\\
 -2&0&0&0&4&0&0&0&0\\
 -2&2&0&0&0&4&0&0&0\\
 -4&0&0&0&0&0&4&0&0\\
 -2&4&0&0&0&0&0&4&0\\
 -4&2&0&0&0&0&0&0&4
 \end{pmatrix}.
 \tag{E67.22}
\]
The independent exact Smith calculation has integer
diagonal entries
\[
 1,1,63,63,63,63,63,2079,2079.
\]
Their three-primary valuations are
$0,0,2,2,2,2,2,3,3$, proving (E67.8).
The supplied certificate retains an integral
three-adic left and right inverse transformation for
the same matrix.

More economically, the explicit character
\[
 \ell=e_{A_0}^*+e_{A_1}^*-e_{B_2}^*\pmod3
       :B^{\rm ab}\longrightarrow\mathbb F_3
 \tag{E67.23}
\]
kills $(16F_B-1)B^{\rm ab}$. On the two columns
of the matrix $U$ in E.37 its value is $(0,-3)$.
The actual Frobenius gluing block is $-2U$, on which
its value is $(0,6)$. The remaining columns of
$16F_B-1$ are divisible by three. But
$\mathfrak m_1(c_*)=[b_0]=e_{A_0}$ and
$\ell(e_{A_0})=1$. Hence $c_*$ represents a
nonzero class in $\mathcal Z_Q$.
Its commutator with $d$ lies in the actual subgroup
$\gamma_3I$, so (E67.10) satisfies the literal tame
relation. Formula (E67.7) excludes every full lift,
including all possible conjugate representatives.

The program \texttt{a\_marked\_local.py} checks the
actual quotient matrix, both exact Smith transforms,
the moment and character identities, the class-two
word and tame equation, and the weighted trace
divisibility used above. Every reported computation
is an exact integral or finite-coefficient calculation.
\end{proof}
